% Provenance (never printed): NSI terminale "Structures de donnees: vocabulaire de la programmation objet (classes,
% attributs, methodes, objets)"; "Langages et programmation: paradigmes de programmation". See sources/book2/BRIEFS.md.
\chapter{Classes and Objects}\label{ch:g12:classes-and-objects}

A date is three numbers: a year, a month, a day. But adding one day to 28 February is not adding one to the third
number: it gives 29 February or 1 March, depending on the year, and 31 December gives a new year. A date is three
numbers \emph{and} the rules that go with them. If the numbers travel through a program without their rules, sooner
or later some function forgets one, and the 30th of February appears.

Python lets a program define its own kinds of values, which carry their data and the functions that respect their
rules together. This chapter defines such classes, protects their rules, gives them the operators of the language,
derives one \omterm{def:g12:classes-and-objects:class}{class} from another, and closes on the different styles in which the same program can be written.

\begin{recall}
A module offers an interface and hides its \omterm{def:g11:modules-and-interfaces:interface}{implementation} (chapter \ref{ch:g11:modules-and-interfaces}). A list is
\omterm{def:g10:tuples-sets-comprehensions:mutable}{mutable}: two names can refer to the same list (chapter \ref{ch:g10:tuples-sets-comprehensions}). The town library's
loans table (chapter \ref{ch:g10:data-tables}).
\end{recall}

\section{Grouping data with the functions that act on it}\label{sec:g12:classes-and-objects:classes}

\begin{definition}[Class, instance, attribute, method]\label{def:g12:classes-and-objects:class}
A \emph{class}\index{class} is a new type of value, defined by the program. A value of that type is an
\emph{instance}\index{instance (object)} of the class, also called an object. Each instance holds named values, its
\emph{attributes}\index{attribute}; the functions defined in the class, which act on an instance, are its
\emph{methods}\index{method}.
\end{definition}

\begin{definition}[Constructor]\label{def:g12:classes-and-objects:constructor}
The \emph{constructor}\index{constructor} of a \omterm{def:g12:classes-and-objects:class}{class} is the \omterm{def:g12:classes-and-objects:class}{method} that gives a new \omterm{def:g12:classes-and-objects:class}{instance} its first \omterm{def:g12:classes-and-objects:class}{attributes}.
In Python it is named \code{__init__}, and it runs when the \omterm{def:g12:classes-and-objects:class}{class} is called: \code{Point(3, 4)}.
\end{definition}

A \omterm{def:g12:classes-and-objects:class}{method} receives the \omterm{def:g12:classes-and-objects:class}{instance} it acts on as its first parameter, by convention named \code{self}:
\code{a.distance_to(b)} is the call \code{Point.distance_to(a, b)}.

\omcode{code/grade-12/02-classes-and-objects/listings/point_class.py}{1}{26}{A \omterm{def:g12:classes-and-objects:class}{class} of points, three \omterm{def:g12:classes-and-objects:class}{instances}, and
two names for one object.}\label{lst:g12:classes-and-objects:point}

\omoutput{code/grade-12/02-classes-and-objects/outputs/point_class.txt}{Moving \code{c} moved \code{b}: they are one
object. \code{d} has the same coordinates, but is another object.}\label{out:g12:classes-and-objects:point}

\begin{omfigure}
\begin{tikzpicture}[font=\footnotesize, box/.style={draw, rounded corners=2pt, align=left, inner sep=4pt}]
  \node[box, fill=omDef!10] (cls) at (0,0) {\textbf{class} \code{Point}\\[2pt] \code{__init__(self, x, y)}\\
    \code{distance_to(self, other)}\\ \code{move(self, dx, dy)}};
  \node[box, fill=black!5] (a) at (5.2,1.3) {\code{x = 0}\\ \code{y = 0}};
  \node[box, fill=black!5] (b) at (5.2,-0.1) {\code{x = 4}\\ \code{y = 5}};
  \node[box, fill=black!5] (d) at (5.2,-1.5) {\code{x = 4}\\ \code{y = 5}};
  \foreach \o in {a, b, d} \draw[->, dashed, black!50] (\o.west) -- (cls.east);
  \node[left] (na) at (8.2,1.3) {\code{a}}; \draw[->] (na) -- (a);
  \node[left] (nb) at (8.2,0.1) {\code{b}}; \node[left] (nc) at (8.2,-0.4) {\code{c}};
  \draw[->] (nb) -- (b.east); \draw[->] (nc) -- (b.east);
  \node[left] (nd) at (8.2,-1.5) {\code{d}}; \draw[->] (nd) -- (d);
  \node[font=\scriptsize, black!60] at (2.9,-2.3) {dashed: ``is an instance of''};
\end{tikzpicture}

{\small One \omterm{def:g12:classes-and-objects:class}{class}, three objects. The \omterm{def:g12:classes-and-objects:class}{class} holds the \omterm{def:g12:classes-and-objects:class}{methods} once; each object holds its own \omterm{def:g12:classes-and-objects:class}{attributes}. Two names,
\code{b} and \code{c}, refer to the same object; \code{d} is a different object with equal values.}
\end{omfigure}

\section{Encapsulation and the interface of a class}\label{sec:g12:classes-and-objects:encapsulation}

\begin{definition}[Encapsulation, class invariant]\label{def:g12:classes-and-objects:encapsulation}
\emph{Encapsulation}\index{encapsulation} is keeping an object's \omterm{def:g12:classes-and-objects:class}{attributes} behind its \omterm{def:g12:classes-and-objects:class}{methods}: users call the
\omterm{def:g12:classes-and-objects:class}{methods}, which are the \omterm{def:g12:classes-and-objects:class}{class}'s interface, and leave the \omterm{def:g12:classes-and-objects:class}{attributes} alone. A \emph{class
invariant}\index{class invariant} is a property of the \omterm{def:g12:classes-and-objects:class}{attributes} that the \omterm{def:g12:classes-and-objects:constructor}{constructor} establishes and every \omterm{def:g12:classes-and-objects:class}{method}
keeps.
\end{definition}

A fraction is kept reduced, with a positive denominator: $6/8$ is stored as $3/4$ and $1/(-3)$ as $-1/3$. The
\omterm{def:g12:classes-and-objects:constructor}{constructor} reduces by the greatest common divisor (the mathematics book), and every operation builds its result
through the \omterm{def:g12:classes-and-objects:constructor}{constructor}, so the invariant holds for every fraction ever made. The \omterm{def:g12:classes-and-objects:class}{attributes} start with an
underscore: they are internal, as in a module (chapter \ref{ch:g11:modules-and-interfaces}).

\omcode{code/grade-12/02-classes-and-objects/listings/fraction.py}{1}{20}{Fractions: the \omterm{def:g12:classes-and-objects:constructor}{constructor} establishes the
invariant.}\label{lst:g12:classes-and-objects:fraction}

The rest of the \omterm{def:g12:classes-and-objects:class}{class} is a set of \omterm{def:g12:classes-and-objects:special-method}{special methods}, explained in the next section.

\omcode{code/grade-12/02-classes-and-objects/listings/fraction.py}{21}{41}{Fractions: printing, comparing, adding,
multiplying.}\label{lst:g12:classes-and-objects:fraction-ops}

\omcode{code/grade-12/02-classes-and-objects/listings/use_fraction.py}{1}{11}{Using the \omterm{def:g12:classes-and-objects:class}{class}: two fractions, their
sum and product, a sort, a conversion, and a denominator $0$.}\label{lst:g12:classes-and-objects:use-fraction}

\omoutput{code/grade-12/02-classes-and-objects/outputs/use_fraction.txt}{Every result is reduced; a denominator $0$ is
refused at once.}\label{out:g12:classes-and-objects:fraction}

The \omterm{def:g12:classes-and-objects:class}{class} \code{Date} keeps a stronger invariant: its three numbers always form a real date. The \omterm{def:g12:classes-and-objects:constructor}{constructor} refuses
anything else, and \code{next_day} handles the month and year changes in one place.

\omcode{code/grade-12/02-classes-and-objects/listings/date_class.py}{1}{35}{Dates that are always
valid.}\label{lst:g12:classes-and-objects:date}

\omcode{code/grade-12/02-classes-and-objects/listings/use_date.py}{1}{10}{Using \code{Date}: four days after,
two comparisons, three impossible dates.}\label{lst:g12:classes-and-objects:use-date}

\omoutput{code/grade-12/02-classes-and-objects/outputs/use_date.txt}{The day after four dates, two comparisons, and
three dates refused.}\label{out:g12:classes-and-objects:date}

Python does not lock the \omterm{def:g12:classes-and-objects:class}{attributes}. A program that writes into \code{_num} breaks the invariant, and with it the
\omterm{def:g12:classes-and-objects:class}{methods} that rely on it.

\omcode{code/grade-12/02-classes-and-objects/listings/bad_mutation.py}{1}{6}{Reaching inside an
object.}\label{lst:g12:classes-and-objects:bad}

\omoutput{code/grade-12/02-classes-and-objects/outputs/bad_mutation.txt}{A ``fraction'' $2/2$, not equal to $1$:
\code{__eq__} assumed the invariant.}\label{out:g12:classes-and-objects:bad}

\begin{remark}
Python also offers \emph{properties}: an \omterm{def:g12:classes-and-objects:class}{attribute} that is read like \code{f.numerator} but computed by a \omterm{def:g12:classes-and-objects:class}{method},
which can refuse to be written. They let a \omterm{def:g12:classes-and-objects:class}{class} change how it stores its data without changing how users read it.
\end{remark}

\section{Special methods}\label{sec:g12:classes-and-objects:special}

\begin{definition}[Special method]\label{def:g12:classes-and-objects:special-method}
A \emph{special method}\index{special method} is a \omterm{def:g12:classes-and-objects:class}{method} with a name of the form \code{__name__} that Python calls
by itself for an operation of the language: \code{__add__} for \code{+}, \code{__eq__} for \code{==}, \code{__lt__}
for \code{<}, \code{__len__} for \code{len}, \code{__str__} for \code{str} and \code{print}, \code{__repr__} for the
form shown in lists and for programmers.
\end{definition}

\omcode{code/grade-12/02-classes-and-objects/listings/special_methods.py}{1}{35}{Durations that print, compare, sort
and add.}\label{lst:g12:classes-and-objects:special}

\omoutput{code/grade-12/02-classes-and-objects/outputs/special_methods.txt}{A list shows \code{repr}, \code{print}
shows \code{str}; \code{sorted} and \code{max} use \code{<}; \code{sum} uses \code{+}.}%
\label{out:g12:classes-and-objects:special}

\section{Inheritance, a first look}\label{sec:g12:classes-and-objects:inheritance}

\begin{definition}[Inheritance, subclass, overriding]\label{def:g12:classes-and-objects:inheritance}
A \emph{subclass}\index{subclass} is a \omterm{def:g12:classes-and-objects:class}{class} defined from another, its base \omterm{def:g12:classes-and-objects:class}{class}, by \emph{inheritance}\index{inheritance}:
it has all the base \omterm{def:g12:classes-and-objects:class}{class}'s \omterm{def:g12:classes-and-objects:class}{methods}, and may add others. A \omterm{def:g12:classes-and-objects:class}{method} of the subclass with the same name as one of the
base \omterm{def:g12:classes-and-objects:class}{class} \emph{overrides}\index{overriding} it: it replaces it for the \omterm{def:g12:classes-and-objects:class}{instances} of the subclass.
\end{definition}

\omcode{code/grade-12/02-classes-and-objects/listings/shapes.py}{1}{34}{A base \omterm{def:g12:classes-and-objects:class}{class}, three subclasses, and
\code{super()}.}\label{lst:g12:classes-and-objects:shapes}

\omoutput{code/grade-12/02-classes-and-objects/outputs/shapes.txt}{\code{describe} is written once and calls the
right \code{area} for each shape; a square is a rectangle.}\label{out:g12:classes-and-objects:shapes}

\begin{omfigure}
\begin{tikzpicture}[font=\footnotesize, box/.style={draw, rounded corners=2pt, align=left, fill=omDef!8,
    inner sep=4pt}, -{Stealth[open, length=3mm]}]
  \node[box] (s) at (0,0) {\textbf{Shape}\\ \code{area()}: raises\\ \code{describe()}};
  \node[box] (r) at (-2.6,-2.2) {\textbf{Rectangle}\\ \code{__init__(w, h)}\\ \code{area()} \textcolor{omThm}{overrides}};
  \node[box] (c) at (2.6,-2.2) {\textbf{Circle}\\ \code{__init__(r)}\\ \code{area()} \textcolor{omThm}{overrides}};
  \node[box] (q) at (-2.6,-4.4) {\textbf{Square}\\ \code{__init__(side)}: \code{super()}};
  \draw (r) -- (s); \draw (c) -- (s); \draw (q) -- (r);
\end{tikzpicture}

{\small The classes of \code{shapes.py}. An arrow goes from a \omterm{def:g12:classes-and-objects:inheritance}{subclass} to its base \omterm{def:g12:classes-and-objects:class}{class}. \code{describe} is
inherited by all; \code{area} is overridden; \code{Square} only changes the \omterm{def:g12:classes-and-objects:constructor}{constructor}, through \code{super()}.}
\end{omfigure}

\omterm{def:g12:classes-and-objects:inheritance}{Inheritance} suits an ``is a'' relation: a square is a rectangle. It is a poor tool for sharing code between classes
that are not of the same kind; then an object that holds another (a playlist holds durations) is clearer.

\section{Programming paradigms}\label{sec:g12:classes-and-objects:paradigms}

\begin{definition}[Programming paradigm]\label{def:g12:classes-and-objects:paradigm}
A \emph{programming paradigm}\index{programming paradigm} is a style of writing programs. In the \emph{imperative}
style, a program is a sequence of instructions that change variables. In the \emph{functional} style, it is built
from functions applied to values, without changing variables. In the \emph{object} style, it is organised around
objects that hold data and the \omterm{def:g12:classes-and-objects:class}{methods} acting on it.
\end{definition}

\begin{definition}[Pure function]\label{def:g12:classes-and-objects:pure-function}
A \emph{pure function}\index{pure function} is a function whose result depends only on its arguments and which
changes nothing outside itself: no global variable, no \omterm{def:g12:classes-and-objects:class}{attribute}, no printing, no file.
\end{definition}

The town library wants the total number of days late over the loans returned after their due date. The same program
three times:

\omcode{code/grade-12/02-classes-and-objects/listings/paradigms.py}{1}{40}{Days late, imperative, functional and
object.}\label{lst:g12:classes-and-objects:paradigms}

\omoutput{code/grade-12/02-classes-and-objects/outputs/paradigms.txt}{Three styles, one
answer.}\label{out:g12:classes-and-objects:paradigms}

\begin{omfigure}
\begin{tikzpicture}[font=\footnotesize, col/.style={draw, rounded corners=2pt, text width=40mm, align=flush left,
    inner sep=4pt, minimum height=27mm, anchor=north}]
  \node[col, fill=omDef!8] at (0,0) {\textbf{imperative}\\[2pt] \code{total = 0}\\ for each loan:\\ \quad if late:
    add its days\\[2pt] \emph{state changed in steps}};
  \node[col, fill=omMeth!10] at (4.6,0) {\textbf{functional}\\[2pt] \code{filter(is_late, LOANS)}\\
    \code{map(days late, ...)}\\ \code{sum(...)}\\[2pt] \emph{whole sequences}};
  \node[col, fill=omThm!8] at (9.2,0) {\textbf{object}\\[2pt] class \code{Loan}\\ method \code{days_late()}\\
    \code{sum(...)} over the loans\\[2pt] \emph{objects answer}};
\end{tikzpicture}

{\small The same computation in three styles. Real programs mix them: the object version uses a functional
\code{sum}, and the \omterm{def:g12:classes-and-objects:class}{methods} of a \omterm{def:g12:classes-and-objects:class}{class} are often written imperatively.}
\end{omfigure}

Each style has its field. The functional style makes programs easy to test and to run in parallel, since \omterm{def:g12:classes-and-objects:pure-function}{pure
functions} cannot disturb one another. The object style suits programs with many kinds of things, each with its own
rules: windows and buttons, accounts, the pieces of a game. The imperative style is the closest to the machine, and
often the most direct for a short computation.

\section{Exercises}\label{sec:g12:classes-and-objects:exercises}

\begin{exercise}[$\star$]\label{exo:g12:classes-and-objects:1}
What does \code{print(Point(1, 2))} show, with the \omterm{def:g12:classes-and-objects:class}{class} of the chapter? How can it be made to show \code{(1, 2)}?
\end{exercise}

\begin{exercise}[$\star$]\label{exo:g12:classes-and-objects:2}
After \code{a = Fraction(1, 2)}, \code{b = a + Fraction(1, 4)} and \code{c = b * Fraction(2, 3)}, what do
\code{print(b, c)} and \code{print(repr(c))} show?
\end{exercise}

\begin{exercise}[$\star$]\label{exo:g12:classes-and-objects:3}
Write a \omterm{def:g12:classes-and-objects:class}{class} \code{Tally} with a \omterm{def:g12:classes-and-objects:class}{method} \code{add(k=1)}, which adds \code{k} to a count starting at $0$, and a \omterm{def:g12:classes-and-objects:class}{method}
\code{value()}.
\end{exercise}

\begin{exercise}[$\star$]\label{exo:g12:classes-and-objects:4}
\code{p = Point(1, 2)}, \code{q = p}, \code{q.move(2, 0)}, \code{r = Point(p.x, p.y)}, \code{r.move(0, 1)}. Give the
coordinates of \code{p}, \code{q} and \code{r}.
\end{exercise}

\begin{exercise}[$\star$]\label{exo:g12:classes-and-objects:5}
A student writes \code{def move(dx, dy):} in the \omterm{def:g12:classes-and-objects:class}{class} \code{Point}. The call \code{p.move(1, 1)} raises
\code{TypeError: Point.move() takes 2 positional arguments but 3 were given}. Explain the three.
\end{exercise}

\begin{exercise}[$\star\star$]\label{exo:g12:classes-and-objects:6}
Write \code{sub(a, b)}, the difference of two fractions, using only the interface of \code{Fraction}. Why is the result
reduced without any work of yours?
\end{exercise}

\begin{exercise}[$\star\star$]\label{exo:g12:classes-and-objects:7}
\code{{Fraction(1, 2)}} raises \texttt{TypeError: unhashable type: \textquotesingle{}Fraction\textquotesingle{}}. Why can an object with \code{__eq__}
not go into a set, until the \omterm{def:g12:classes-and-objects:class}{class} also defines \code{__hash__}?
\end{exercise}

\begin{exercise}[$\star\star$]\label{exo:g12:classes-and-objects:8}
\code{A.__init__} prints \code{A}. \code{B(A)} has an \code{__init__} that prints \code{B in}, calls
\code{super().__init__()}, then prints \code{B out}. What does \code{B()} print? And without the call to
\code{super()}?
\end{exercise}

\begin{exercise}[$\star\star$]\label{exo:g12:classes-and-objects:9}
Write a \omterm{def:g12:classes-and-objects:class}{method} \code{previous_day} for the \omterm{def:g12:classes-and-objects:class}{class} \code{Date}, and check it on 1 March 2028 and 1 January 2027.
\end{exercise}

\begin{exercise}[$\star\star$]\label{exo:g12:classes-and-objects:10}
Prove that every \code{Fraction}, including the results of \code{+} and \code{*}, has a positive denominator and a
numerator and denominator with no common divisor other than $1$.
\end{exercise}

\begin{exercise}[$\star\star$]\label{exo:g12:classes-and-objects:11}
Rewrite in the functional style, with \code{filter}, \code{map} and \code{sum}: \code{total = 0}, then
\code{for x in t:}, \texttt{if x \% 2 == 0:}, \code{total += x * x}.
\end{exercise}

\begin{exercise}[$\star\star\star$]\label{exo:g12:classes-and-objects:12}
Write a \omterm{def:g12:classes-and-objects:class}{class} \code{Stack} with \code{push}, \code{pop}, \code{peek}, \code{is_empty} and \code{len}, where
\code{pop} and \code{peek} raise \code{IndexError} on an empty stack, and tests for it.
\end{exercise}

\begin{exercise}[$\star\star\star$]\label{exo:g12:classes-and-objects:13}
How many \code{Fraction} objects does this program create? \code{total = Fraction(0)}, then for $k$ from $1$ to $4$,
\code{total = total + Fraction(1, k)}. How many are still reachable at the end?
\end{exercise}

\begin{exercise}[$\star\star\star$]\label{exo:g12:classes-and-objects:14}
A \omterm{def:g12:classes-and-objects:class}{class} \code{Basket} has the line \code{items = []} in its body, outside any \omterm{def:g12:classes-and-objects:class}{method}, and a \omterm{def:g12:classes-and-objects:class}{method} \code{add} that
appends to \code{self.items}. After \code{a.add("pear")}, the basket \code{b} also holds a pear. Explain and correct.
\end{exercise}

\begin{exercise}[$\star\star\star$]\label{exo:g12:classes-and-objects:15}
Write a \omterm{def:g12:classes-and-objects:class}{class} \code{Polynomial} built from its list of coefficients $[a_0, a_1, \dots]$, with \code{__add__},
\code{__mul__}, and \code{__call__} so that \code{p(x)} computes the value at \code{x}.
\end{exercise}

\section{Problem: Exact Fractions}\label{sec:g12:classes-and-objects:problem}

\begin{problem}[{Weekend problem --- the harmonic sums, exactly}]\label{pb:g12:classes-and-objects:1}
The harmonic sums are $H_n = 1 + \frac12 + \frac13 + \dots + \frac1n$. With the \omterm{def:g12:classes-and-objects:class}{class} \code{Fraction}, a program
computes them exactly, compares their denominators with the least common multiple of $1, \dots, n$, and measures the
error of the same sums in floating point.

\omcode{code/grade-12/02-classes-and-objects/listings/harmonic.py}{1}{17}{Harmonic sums exact and in floating
point.}\label{lst:g12:classes-and-objects:harmonic}

\omoutput{code/grade-12/02-classes-and-objects/outputs/harmonic.txt}{The exact sums, their denominators, and three
comparisons with floating point.}\label{out:g12:classes-and-objects:harmonic}

\needspace{4\baselineskip}\medskip\noindent\textbf{Part I --- A class of fractions.}
\begin{enumerate}
  \item Which \omterm{def:g12:classes-and-objects:class}{attributes} does a \code{Fraction} have, and why do their names start with an underscore?
  \item What does \code{Fraction(10, -4)} store?
  \item Which \omterm{def:g12:classes-and-objects:special-method}{special method} does \code{exact + Fraction(1, n)} call, and with which arguments?
  \item Why does \code{Fraction(1, 2) == Fraction(2, 4)} hold, although $2/4$ was written?
  \item Why does the \omterm{def:g12:classes-and-objects:class}{class} define \code{numerator()} and \code{denominator()} rather than letting users read
    \code{_num} and \code{_den}?
\end{enumerate}

\needspace{4\baselineskip}\medskip\noindent\textbf{Part II --- The invariant.}
\begin{enumerate}[resume]
  \item State the invariant of \code{Fraction}.
  \item Prove that the \omterm{def:g12:classes-and-objects:constructor}{constructor} establishes it, whatever the signs of its arguments (the denominator not $0$).
  \item Show that \code{__eq__} is right because of the invariant: two fractions are equal numbers exactly when their
    reduced forms are equal.
  \item What does \code{bad_mutation.py} show about Python's protection of \omterm{def:g12:classes-and-objects:class}{attributes}?
\end{enumerate}

\needspace{4\baselineskip}\medskip\noindent\textbf{Part III --- The harmonic sums exactly.}
\begin{enumerate}[resume]
  \item Compute $H_4$ by hand, and check the output.
  \item From the output, give $H_9$ and $H_{10}$. Check $H_{10} = H_9 + \frac{1}{10}$.
  \item How many digits has the numerator of $H_{1000}$? Why could no type of fixed size hold it?
  \item Why can no \omterm{def:g10:floating-point:float}{floating-point number} hold $H_3 = 11/6$ exactly?
\end{enumerate}

\needspace{4\baselineskip}\medskip\noindent\textbf{Part IV --- Denominators and floating point.}
\begin{enumerate}[resume]
  \item Why is $\operatorname{lcm}(1, \dots, n)$ always a common denominator of the terms of $H_n$?
  \item For which $n$ of the output is the reduced denominator smaller than that common multiple? By how much?
  \item For $n = 6$, write $H_6$ over $60$ and explain why it reduces by $3$.
  \item At $n = 10$ and $n = 100$, how large is the floating-point error printed? And at $n = 1000$?
  \item Where does that error come from? Recall chapter \ref{ch:g10:floating-point}.
  \item Why is the error at $n = 1000$ so small compared with the value?
  \item State the result: $H_{10}$ exactly, and the error of the floating-point sum up to $1000$.
\end{enumerate}
\end{problem}
